\section{Experimental Evaluation}
\label{sec:eval}
The evaluation asks whether one registered-harness path can execute bounded
tasks across two simulator scenes with different signal and protocol mappings
while retaining a checkable task-to-approval-to-receipt-to-sample trace. R20 is
the primary result: six sequential tasks through one registered harness, three
per scene on one simulator instance per scene. It demonstrates this configured
software path, not configuration-reuse benefit or process attainment. R21 is a
separate deterministic fixture replay on an isolated current-version
application and simulator; it checks the same task-to-device path without an
LLM decision. R13 is secondary exploratory protocol and data-path evidence
whose exact runtime source was not pinned. All runs use repository simulators.
Supplementary File~S1 contains the R20 trace matrix and records, plus the R21
fixture evidence.

\subsection{Evidence Units and Experimental Procedure}
\label{sec:testbed}
\label{sec:procedure}
R20 uses the registered harness adapter named \texttt{omp} in cast-film and annealing scenes
(\texttt{omp} is the project adapter key, not a network protocol),
with three sequential task executions per scene on one simulator instance per
scene. Each execution creates a new logical line, product, recipe, active run,
channel, agent, and task on that shared scene simulator. Line start reapplies the recipe baseline
(120~rpm or 730~$^{\circ}$C), so the control read returns to baseline while
the simulated process state carries forward. These are sequential software
runs, not independent process replications. R13 uses configured protocol and
data paths in the same two simulator scenes. Earlier experiments are retained
separately in Supplementary File~S1 and are not pooled with these results.

\subsection{Registered-Harness Process Path (R20)}
\label{sec:harness-e2e}
R20 ran the local v0.8.0 working-tree application with the registered
\texttt{omp} adapter and GLM-5.3-Flash. Supplementary File~S1 contains the
source-reconstruction kit and lists its external prerequisites: the base
repositories, OMP CLI, local model-service credentials, Node.js toolchain, and
simulator runtime. Credentials are not included, and the external model
service does not guarantee deterministic replay. Cast film used OPC UA for screw-speed
control and MQTT for thickness; annealing used Modbus TCP for both the zone-3
setpoint and strip-temperature probe. The model received a live control read,
process measurement, and target band, then selected one value from the
scenario-defined prevalidated set (125/130/135 rpm or 735/740/745~$^{\circ}$C)
for a bounded \texttt{recipe\_trial}; this was a constrained choice, not an
unrestricted optimization search. A synthetic administrator approved each
proposal. Each route log contains one trial and a 1/1 verified device receipt.
Workspace history can mirror host-tool events; trial counts therefore use the
route-specific recipe and acknowledgment records.

After verifying the route ACK and returned trial result, the runner initiated
a 20-s delay, then probed, ingested, and queried one value per task under the
active-run tag through the isolated SQLite time-series fallback. The runner
checked task state after the sample query; all six task records were then
\texttt{COMPLETED}. The runner reported 76 passing assertions and no failures, and
the separate trace audit found a complete task-to-sample record in all six
executions: each had a matching approval ID, one readback-verified trial
receipt, and one query result whose run ID and value matched the active run and
explicit probe. The task-completion timestamp was not retained, so no
task-completion-to-sample interval can be computed. Table~\ref{tab:q8} labels
Film-1--3 map to \texttt{cast-film-r1--r3}, and Anneal-1--3 to
\texttt{annealing-r1--r3}; full identifiers and per-record checks are listed
in \texttt{R20-TRACE-MATRIX.csv} in S1.

The trace matrix distinguishes two run identifiers. The trial run ID denotes
the \texttt{recipe\_trial} operation record; the active run ID denotes the
production run used to tag and query the probe. Both are linked to the same
task and approval record, and they are not expected to be equal.

Table~\ref{tab:q8} reports the pre-task process value, selected setting, and
run-tagged sample for each execution; all six remained outside target. The
results verify a bounded harness-to-simulator path, not process attainment or
convergence. One probe is not a response trajectory, the SQLite fallback does
not test normal online DAQ, and R20 does not test profile/binding reuse or
cross-harness portability.

\begin{table}[!t]
\centering
\caption{Constrained model selections and simulator measurements. The control
is screw speed for cast film and the zone-3 furnace setpoint for annealing.
Each row reports one run-tagged probe after a 20-s delay from verified route
readback. All six samples are outside target; executions are sequential within
each scene.}
\label{tab:q8}
\footnotesize
\setlength{\tabcolsep}{2.5pt}
\begin{tabularx}{\columnwidth}{@{}lXcX@{}}
\toprule
Case & Pre-task value / target band & Control: start$\to$trial & Sample \\
\midrule
Film-1 & $55.82/48$--$52~\mu$m & $120\to125$~rpm & $45.69~\mu$m \\
Film-2 & $45.69/48$--$52~\mu$m & $120\to130$~rpm & $46.69~\mu$m \\
Film-3 & $46.69/48$--$52~\mu$m & $120\to125$~rpm & $45.69~\mu$m \\
Anneal-1 & $675.40/700$--$760~^{\circ}$C & $730\to745$~$^{\circ}$C & $679.60~^{\circ}$C \\
Anneal-2 & $678.30/700$--$760~^{\circ}$C & $730\to745$~$^{\circ}$C & $677.70~^{\circ}$C \\
Anneal-3 & $676.90/700$--$760~^{\circ}$C & $730\to745$~$^{\circ}$C & $678.10~^{\circ}$C \\
\bottomrule
\end{tabularx}
\end{table}

\subsection{Deterministic Fixture Replay (R21)}
R21 ran the local v0.8.0 application with a deterministic transport fixture
for the registered \texttt{omp} harness,
which returned a fixed tool call without invoking an LLM or autonomous harness
decision. In one fresh isolated application/simulator environment, one fixture
task per scene invoked the recipe-trial path in cast-film and annealing. The
runner checked dispatch, synthetic approval, verified receipt, matching device
readback, run-tagged DAQ probe, and route log; all 31 runner assertions,
including process and port cleanup checks, passed. This replay shows that the
configured software route can be exercised with deterministic harness input.
It is not an independent real-agent run, a process replication, or evidence of
target attainment. The two fixture tasks are not pooled with R20.

\subsection{Software Contract Checks}
The local unit suite passed 554/554 tests across seven suites. Relevant
contract checks refuse denied or expired approvals, changed recipe or limit
state, and invalid node values before driver actuation; unchanged approved
state dispatches through the mock driver and readback. The approval-snapshot
suite passed 20/20 cases. Twelve deterministic fault injections changed one
approval-time field at a time (line, product, recipe, or active run) on
\texttt{recipe\_apply}, \texttt{recipe\_trial}, and
\texttt{recipe\_rollback(dispatch=true)}. Every stale action was refused
before a mock-driver write or new execution record. The handler compares the
saved line--product--recipe--run tuple after the approval wait. These checks
use temporary or mock state and are separate from R20's 76 runner checks;
they do not represent process trials or test missing/mismatched DAQ run tags
in the end-to-end path.

\subsection{Exploratory Protocol and Data Path (R13)}
\label{sec:crossscene}
The R13 application report describes five protocol-family read probes (OPC UA,
MQTT, Modbus TCP, Modbus RTU over TCP, and HTTP/REST) and bounded writes with
matching readbacks in two simulator scenes. Six sequential 120-s windows per
scene at 1-s sampling produced lineage-verified datasets with 658 cast-film
rows and 655 annealing rows. The rows are overlapping time-window samples, not
independent process replications. The exact runtime source was not recorded. R13 is
used only as exploratory protocol and data-lineage evidence; its archived gate
outputs and run-boundary audit are documented in the reproducibility
supplement and do not support the primary registered-harness result.
R14 is an earlier exploratory pilot with alternative signal pairs; both pilot
model gates failed, and those outcomes informed the subsequent R13 pair
selection. R13 also produced no eligible model, recommendation, or dispatch.
The two runs remain separate and are not model-performance evidence.

Earlier protocol, admission-ablation, controller, deployment, and scripted
recipe records in S1 use older versions or interfaces. They are not pooled
with R13 or R20, used as matched strategy comparisons, or relied on for the
current integration claims.
